\section{The instrument: gTAO-Net}
\label{sec:instrument}

gTAO-Net (gated Tchebichef Attention Ordinal network) implements each design choice as a switch. For a window $x=\{x_k\}_{k=1}^{S}$ with $S$ streams, the logits are
\begin{equation}
  \ell(x)=h\big(F(T(x))+G(x)\big),
\end{equation}
where $T$ denotes the stream towers, $F$ the fusion operator, $G$ an optional global link and $h$ the ordinal head (Fig.~\ref{fig:arch}).

\begin{figure}[!htbp]
\centering
\resizebox{0.85\linewidth}{!}{%
\begin{tikzpicture}[
  font=\sffamily\small, >={Stealth[length=2.2mm]},
  box/.style={draw, rounded corners=2pt, minimum height=7mm, align=center, inner sep=2.5pt},
  inp/.style={box, fill=gray!8, minimum width=22mm},
  mom/.style={box, fill=orange!15, minimum width=25mm},
  tow/.style={box, fill=blue!8, minimum width=15mm},
  op/.style={box, fill=violet!10},
  ]
  \foreach \n/\lab/\enc/\dim [count=\i] in
    {face/{Face\\$64{\times}64$}/ {2D16+1D6}/1536,
     fm/{Face mesh\\$478{\times}3$}/{1D6}/8604,
     hp/{Head pose\\3 ang}/{1D6}/18,
     ph/{Physiology\\8 ch}/{1D6}/48} {
    \node[inp] (in\n) at (0,-1.2*\i) {\lab};
    \node[mom, right=5mm of in\n] (m\n) {Tcheb.\\\enc};
    \node[tow, right=7mm of m\n] (t\n) {tower\\$z_{\n}$};
    \draw[->] (in\n) -- (m\n);
    \draw[->] (m\n) -- node[above, font=\sffamily\scriptsize]{\dim} (t\n);
  }
  \node[op, minimum height=40mm, minimum width=28mm, right=9mm of $(tfm)!0.5!(thp)$] (fuse)
    {self-attention\\$Z'=\mathrm{softmax}\!\big(\tfrac{QK^\top}{\sqrt d}\big)Z$\\[2mm]
     \textbf{gate}: $\tilde z=\sum_i \alpha_i Z'_i$\\[1mm] \emph{or}\\[1mm]
     \textbf{core}: $\tilde z=\tfrac1S\sum_i Z'_i$};
  \foreach \n in {face,fm,hp,ph} \draw[->] (t\n.east) -- (fuse.west |- t\n.east);
  \node[draw, circle, inner sep=1pt, right=7mm of fuse] (plus) {$+$};
  \node[box, fill=yellow!20, right=6mm of plus] (head) {CORAL head};
  \node[box, fill=green!10, right=6mm of head] (out) {level};
  \draw[->] (fuse) -- (plus); \draw[->] (plus) -- (head); \draw[->] (head) -- (out);
  \coordinate (lk) at ($(inph.south)+(0,-6mm)$);
  \foreach \n in {face,fm,hp,ph} {
    \coordinate (tap\n) at ($(m\n.east)!0.72!(t\n.west)$);
    \fill[red!70!black] (tap\n) circle (1.2pt);
  }
  \draw[very thick, red!70!black] (tapface) -- (tapph);
  \draw[->, very thick, red!70!black] (tapph) -- (tapph |- lk) -| node[pos=0.25, below, font=\sffamily\scriptsize] {global link $G(x)=\sum_k W^{g}_k\phi_k(x_k)$} (plus.south);
  \node[below=1mm of head, font=\sffamily\scriptsize] {CORAL $|$ hybrid $|$ soft};
\end{tikzpicture}}
\caption{gTAO-Net: fixed Tchebichef front-ends $\phi_k$, stream towers, gate or core fusion, an optional linear global link, and a CORAL head trained with one of three objectives.}
\label{fig:arch}
\end{figure}

\textbf{Front-ends and towers.} Each stream is a short per-frame sequence. Face-stream frames are converted to grey $64\times64$ images and projected on the 2-D Tchebichef basis of order 16~\cite{mukundan2001tchebichef}; every coefficient, landmark coordinate, angle and physiological channel is then projected in time on the 1-D basis of order 6. This yields fixed vectors of 1{,}536 (face), 8{,}604 (mesh), 18 (pose) and 48 (physiology) dimensions, which are computed once and cached; no front-end parameter is learned. One transform family thus serves all four modalities, so adding a stream requires choosing a moment order rather than training an encoder~\cite{makhkhas2026ordinal}. Each tower computes $z_k=\mathrm{LN}(W_2\,\mathrm{GELU}(W_1\phi_k))\in\mathbb{R}^{64}$ (hidden size 128). Single-stream baselines consist of one tower followed by the head.

\textbf{Fusion.} The stacked summaries $Z$ pass through one self-attention step with bias-free query and key projections, $Z'=\mathrm{softmax}(QK^\top/\sqrt{d})Z$. The \emph{gate} maps $Z'$ through a two-layer MLP $g$ to weights $\alpha=\mathrm{softmax}(g(\mathrm{vec}\,Z'))$ and returns $\tilde z=\sum_i\alpha_iZ'_i$; these weights allow it to discard a stream. The \emph{core} variant, following CORE-Net~\cite{tran2026ordinal}, returns the mean $\tilde z=S^{-1}\sum_iZ'_i$ and cannot discard a stream. Modality dropout~\cite{neverova2016moddrop} zeroes each tower output with probability 0.2 during training, independently per sample and stream and without rescaling; the global link receives the unmasked inputs.

\textbf{Global link and objectives.} The global link $G(x)=\sum_kW^{g}_k\phi_k(x_k)$ is linear and bias-free, so its gradient reaches the input in two steps, as in a residual shortcut~\cite{he2016deep}. The CORAL head~\cite{cao2020coral} computes $\ell_k=w^\top u+b_k$ with $u=\tilde z+G(x)$. Three objectives are compared: \textbf{CORAL}, a binary cross-entropy between $\sigma(\ell_k)$ and $\mathbf{1}[y>k]$; \textbf{hybrid}, inspired by CORE-Net~\cite{tran2026ordinal}, which adds to the cross-entropy on the implied class probabilities 0.5 times the CORAL loss and 0.1 times a cross-entropy to a neighbour-smoothed target (0.7 on the true class and 0.15 on each neighbour, renormalised at the extremes); unlike CORE-Net, it has no earth mover's term and gives some probability to neighbouring classes instead of penalising it; and \textbf{soft}, a binary cross-entropy against $t_k=1-\sum_{j\le k}\bar p_j$.

\noindent\textbf{Proposition.} If $\bar p$ is one-hot at $y$, then $t_k=\mathbf{1}[y>k]$ and the soft objective equals CORAL. Consequently, soft and CORAL are identical on the single-label datasets DAiSEE and EngageNet (confirmed bit-identical in the logs), and on DIPSER any difference between them arises from rater disagreement alone.
